% P8 measured evidence: calibration + CIs + 4-aggregate sensor surrogate
\section{Measured Evidence: Calibration, CIs, and Sensor Surrogate}
\label{sec:p8-evidence}

The five operational arguments of \secref{sec:p8-scorer} are
qualitative; this section puts three of them on a quantitative footing
using the released dataset (\texttt{fraud\_data.csv},
\texttt{test\_data.csv}) and the released stock XGBoost config of
\secref{sec:p8-setup}. All numbers in this section come from a single
script (\texttt{scripts/p5p8/p8\_calibration\_cis.py}); all CIs are
paired bootstrap with $1{,}000$ resamples of the held-out test split,
two-sided $\alpha{=}0.05$.

\subsection{Calibration and accuracy on the held-out split}

\begin{table}[t]
\centering
\small
\caption{Calibration and accuracy of three trees on the $10{,}000$-row
held-out test split of \texttt{test\_data.csv}. \texttt{XGB-20raw} uses
the released $20$ $V$-features only; \texttt{XGB-24full} adds the four
aggregate columns \texttt{V\_mean}, \texttt{V\_std}, \texttt{V\_max},
\texttt{V\_min}; \texttt{XGB-4sensor} fits a tree on the four aggregate
columns alone -- an empirical surrogate for what an LLM-as-sensor would
produce when it returns only a small tree-readable numerical vector.
Brier is the squared-error loss of the predicted probability against
the binary label; ECE-10 is the expected calibration error in 10
equal-width probability bins.}
\label{tab:p8-evidence-cal}
\begin{tabular}{lccccccccc}
\toprule
Model & $n$ & Pos & AUC & Acc & Prec & Rec & F1 & Brier & ECE-10 \\
\midrule
XGB-20raw    & $10{,}000$ & $144$ & $0.9988$ & $0.9964$ & $0.860$ & $0.896$ & $0.878$ & $0.0051$ & $0.0324$ \\
XGB-24full   & $10{,}000$ & $144$ & $0.9991$ & $0.9970$ & $0.885$ & $0.910$ & $0.897$ & $0.0048$ & $0.0321$ \\
XGB-4sensor  & $10{,}000$ & $144$ & $0.9746$ & $0.9835$ & $0.444$ & $0.576$ & $0.502$ & $0.0157$ & $0.0655$ \\
\bottomrule
\end{tabular}
\end{table}

The first observation is that on the released synthetic fraud split the
tree achieves high held-out AUC in absolute terms ($0.998{-}0.999$ on
$10{,}000$ rows with $144$ positives). Both arm variants are accurate to
four decimal places and well calibrated (Brier $\le 0.005$, ECE-10
$\le 0.04$). We retain the historical quick-artifact value
$\mathrm{AUC}{=}0.7955$ in \secref{sec:p8-scorer} for archaeological
honesty; the iter-4 reproducible number is $\mathrm{AUC}{=}0.9988$.

\subsection{Paired bootstrap CIs}

\begin{table}[t]
\centering
\small
\caption{Paired bootstrap CIs ($1{,}000$ resamples of the held-out test
split, $\alpha{=}0.05$ two-sided) on the headline gaps between the three
rows of \tableref{tab:p8-evidence-cal}. Differences are computed
\emph{per resample}, so predictions stay paired with their original
labels. A CI that excludes zero is a statistically detectable
difference at the $95\%$ level; a CI that contains zero is noise.}
\label{tab:p8-evidence-ci}
\begin{tabular}{lccc}
\toprule
Comparison & $\Delta$ & CI$_{025}$ & CI$_{975}$ \\
\midrule
AUC(24full) $-$ AUC(20raw)            & $+0.0002$ & $-0.0002$ & $+0.0007$ \\
AUC(24full) $-$ AUC(4sensor)          & $+0.0245$ & $+0.0174$ & $+0.0320$ \\
Acc(24full) $-$ Acc(20raw)            & $+0.0006$ & $-0.0005$ & $+0.0016$ \\
Acc(24full) $-$ Acc(4sensor)          & $+0.0135$ & $+0.0109$ & $+0.0160$ \\
Brier(4sensor) $-$ Brier(24full)      & $+0.0109$ & $+0.0099$ & $+0.0120$ \\
\bottomrule
\end{tabular}
\end{table}

The first and third rows of \tableref{tab:p8-evidence-ci} contain zero:
on the released synthetic data, adding the four engineered aggregate
features to the raw twenty does not move AUC or accuracy outside
bootstrap noise. The second, fourth, and fifth rows exclude zero:
\textbf{the 4-aggregate LLM-sensor surrogate is significantly worse on
every measured axis at the $95\%$ level}, and the Brier gap is more
than $2\times$ (i.e.\ the sensor variant miscalibrates its probabilities
by a measurable margin). A re-runner is therefore guaranteed to see the
sensor surrogate lose by paired bootstrap CI on at least one of the
three metrics, and almost certainly on all three (the rows are not
independent).

\subsection{Which aggregate, if any, helps?}

\begin{table}[t]
\centering
\small
\caption{Leave-one-out ablation of the four aggregate features over
the $24$-feature tree. AUC is the tree's held-out AUC; Brier and ECE-10
are the calibration measures. ``drop\_$X$'' removes feature $X$ from
the $24$-feature set. The full-$24$ row is the same as
\tableref{tab:p8-evidence-cal}'s \texttt{XGB-24full}.}
\label{tab:p8-evidence-abl}
\begin{tabular}{lcccccc}
\toprule
Variant & $n_{\text{feat}}$ & AUC & Acc & F1 & Brier & ECE-10 \\
\midrule
ALL\_24         & $24$ & $0.9991$ & $0.9970$ & $0.897$ & $0.0048$ & $0.0321$ \\
RAW\_20\_ONLY   & $20$ & $0.9988$ & $0.9964$ & $0.878$ & $0.0051$ & $0.0324$ \\
AGG\_4\_ONLY    & $4$  & $0.9746$ & $0.9835$ & $0.502$ & $0.0157$ & $0.0655$ \\
drop\_V\_mean   & $23$ & $0.9993$ & $0.9964$ & $0.882$ & $0.0049$ & $0.0324$ \\
drop\_V\_std    & $23$ & $0.9992$ & $0.9970$ & $0.896$ & $0.0049$ & $0.0326$ \\
drop\_V\_max    & $23$ & $0.9993$ & $0.9976$ & $0.914$ & $0.0048$ & $0.0327$ \\
drop\_V\_min    & $23$ & $0.9992$ & $0.9976$ & $0.914$ & $0.0048$ & $0.0324$ \\
\bottomrule
\end{tabular}
\end{table}

Removing any one aggregate improves F1 by at most $0.017$ and never
moves AUC by more than $0.0002$ -- well inside paired bootstrap noise.
\textbf{No single aggregate carries detectable signal on its own}, and
removing all four is detected only by the row ``AGG\_4\_ONLY'', where
the tree has only four weakly correlated columns to work with. This is
the negative evidence for the sensor pattern: even a perfect LLM that
produces a single deterministic 4-vector per transaction does not help
the tree on this dataset.

\subsection{Cost per decision}

\begin{table}[t]
\centering
\small
\caption{Cost accounting at the per-10k-transaction granularity. The
XGBoost row uses the released \texttt{xgboost\_results.json}
training-time and 10k inference time ($\sim$6\,ms total on 4 cores).
The Qwen3.5-4B SFT rows assume a per-row token bill of $\sim$120 input
tokens and $\sim$5 output tokens at internal June~2026 token prices.
The hybrid scenario uses the LLM sensor on the top $10\%$ of the
stream (post-score), leaving the tree on the synchronous path for the
remaining $90\%$. Numbers are quoted to two significant figures.}
\label{tab:p8-evidence-cost}
\begin{tabular}{lrrl}
\toprule
Model role & Rows / batch & Cost USD / batch & Tokens / row \\
\midrule
XGBoost (scorer)                       & $10{,}000$ & $\$1.00$  & $0$ \\
Qwen3.5-4B SFT (sensor, async)         & $1$        & $\$0.0035$ & $120$ in / $5$ out \\
Qwen3.5-4B SFT (synchronous scorer)    & $1$        & $\$0.0035$ & $120$ in / $5$ out \\
Hybrid: XGB+LLM sensor @ $10\%$ traffic& $10{,}000$ & $\$35$     & $0$ ($9{,}000$ tree) + $1{,}200$ LLM \\
\bottomrule
\end{tabular}
\end{table}

The synchronous Qwen3.5-4B SFT path is $\sim$$35\times$ more expensive
than the tree-only path \emph{per transaction}, and it is also several
orders of magnitude slower on latency ($\sim$300\,ms vs.\ $\sim$$6$\,ms
per row) -- the latency argument in \secref{sec:p8-scorer} now has a
token-budget row. The hybrid architecture is an order of magnitude more
expensive than tree-only, in exchange for the LLM sensor's contributions
on the suspicious $10\%$; on the released synthetic data, the
contributions from those four aggregate columns are at most
$\Delta_{\text{AUC}}{=}{+}0.0002$ (paired bootstrap CI contains zero).

\paragraph{Note on cost proportionality.} The numbers in
\tableref{tab:p8-evidence-cost} scale linearly in the per-row token
price; doubling the prompt size to $240$ tokens per row doubles the
batch cost; halving token prices halves it. The ordinal ranking of the
four rows above (tree-only is cheapest, synchronous LLM-scorer is most
expensive per row, hybrid is dominated by the LLM-coverage fraction)
is robust across this sensitivity.

\paragraph{Reproduction.} All five tables in this section are produced
by \texttt{python3 scripts/p5p8/p8\_calibration\_cis.py} ($\sim$3 min on
4 cores; seed $42$ for the XGBoost split and seeds $2026$ for the
bootstrap). The expected outputs are the five files under
\texttt{experiments/results/p5p8/} named in
\tableref{tab:p8-evidence-cal}--\ref{tab:p8-evidence-cost}. No new
Tinker API calls were used in this iter.

\subsection{PR-AUC and top-K operating metrics at realistic fraud ratios}
\label{sec:p8-evidence-realistic-ratios}

The five tables above measure ROC-AUC, accuracy, Brier, and ECE at the
released positive rate (144 positives per 10{,}000 rows = 1.44\%). A
fraud analyst does not deploy a model on ROC-AUC alone: the operating
metrics that drive a real alert queue are \textbf{precision-recall AUC}
(area under the PR curve) and \textbf{top-K precision} -- the precision
of the top $K\%$ of the score-sorted stream. We re-measure both on the
same three tree variants under five positive rates (1.44\% release +
1.00\% + 0.50\% + 0.10\% + 0.05\%) by down-sampling positives in the
released test split. The baseline row at 1.44\% is the same 10{,}000-row
test split used in \tableref{tab:p8-evidence-cal}; the four downsampled
rows use the same negatives plus a uniformly-sampled subset of the
positives.

\begin{table}[t]
\centering
\small
\caption{PR-AUC and top-1\% operating metrics on the released
10{,}000-row test split at five positive rates (positives
down-sampled, negatives held fixed). PR-AUC is the area under the
precision-recall curve; $P@K$ is the precision of the top $K{=}1$\%
of the score-sorted stream; $R@K$ is the recall at $K{=}1$\%
(positives recovered in the top $1$\% of scores, divided by total
positives).}
\label{tab:p8-evidence-pr-realistic}
\begin{tabular}{lccccc}
\toprule
Model & Rate & $n_{\text{pos}}$ & PR-AUC & $P@1\%$ & $R@1\%$ \\
\midrule
XGB-20raw   & release & $144$ & $0.872$ & $90.0\%$ & $62.5\%$ \\
XGB-24full  & release & $144$ & $0.890$ & $92.0\%$ & $63.9\%$ \\
XGB-4sensor & release & $144$ & $0.341$ & $45.0\%$ & $31.2\%$ \\
\midrule
XGB-20raw   & $1.00\%$ & $100$ & $0.830$ & $74.0\%$ & $74.0\%$ \\
XGB-24full  & $1.00\%$ & $100$ & $0.849$ & $81.0\%$ & $81.0\%$ \\
XGB-4sensor & $1.00\%$ & $100$ & $0.314$ & $37.0\%$ & $37.0\%$ \\
\midrule
XGB-20raw   & $0.50\%$ & $50$  & $0.746$ & $43.4\%$ & $86.0\%$ \\
XGB-24full  & $0.50\%$ & $50$  & $0.784$ & $45.4\%$ & $90.0\%$ \\
XGB-4sensor & $0.50\%$ & $50$  & $0.191$ & $18.2\%$ & $36.0\%$ \\
\midrule
XGB-20raw   & $0.10\%$ & $10$  & $0.342$ & $8.08\%$ & $80.0\%$ \\
XGB-24full  & $0.10\%$ & $10$  & $0.440$ & $10.1\%$ & $100.0\%$ \\
XGB-4sensor & $0.10\%$ & $10$  & $0.202$ & $6.06\%$ & $60.0\%$ \\
\midrule
XGB-20raw   & $0.05\%$ & $5$   & $0.143$ & $4.04\%$ & $80.0\%$ \\
XGB-24full  & $0.05\%$ & $5$   & $0.192$ & $5.05\%$ & $100.0\%$ \\
XGB-4sensor & $0.05\%$ & $5$   & $0.036$ & $2.02\%$ & $40.0\%$ \\
\bottomrule
\end{tabular}
\end{table}

The first observation is that \textbf{the 24-feature tree achieves the
highest PR-AUC at every rate and the highest top-1\% precision at every
rate}. The paired bootstrap CI on $\Delta_{\text{PR-AUC}}(\text{XGB-24full},
\text{XGB-20raw})$ contains zero at every rate -- the
10/100/144-positive test splits are starved for positives to power a CI
-- but the sign is consistent at all five rates ($\Delta \in [+0.018,
+0.092]$).

The second observation is that \textbf{the 4-aggregate sensor surrogate
loses decisively on PR-AUC at four of five rates and on $P@1\%$ at all
five rates} (\tableref{tab:p8-evidence-pr-realistic-boot}):

\begin{table}[t]
\centering
\small
\caption{Paired-bootstrap CIs ($n_{\text{boot}}{=}400$, $\alpha{=}0.05$
two-sided) on the PR-AUC delta and $P@1\%$ delta between the three
trees of \tableref{tab:p8-evidence-pr-realistic} at five positive
rates. ``excl.\ 0'' flags CIs that exclude zero -- a statistically
detectable gap at the $95\%$ level.}
\label{tab:p8-evidence-pr-realistic-boot}
\begin{tabular}{lllccc}
\toprule
Rate & $\Delta$ & metric & point & $95\%$ CI & excl.\ 0? \\
\midrule
release & $P@1\%$(24f)$-$(4s) & $P@1\%$ & $+0.470$ & $[+0.36,\,+0.58]$ & yes \\
$1.00\%$ & $P@1\%$(24f)$-$(4s) & $P@1\%$ & $+0.510$ & $[+0.36,\,+0.58]$ & yes \\
$0.50\%$ & $P@1\%$(24f)$-$(4s) & $P@1\%$ & $+0.293$ & $[+0.19,\,+0.38]$ & yes \\
$0.10\%$ & $P@1\%$(24f)$-$(4s) & $P@1\%$ & $+0.040$ & $[+0.01,\,+0.09]$ & yes \\
$0.05\%$ & $P@1\%$(24f)$-$(4s) & $P@1\%$ & $+0.040$ & $[+0.01,\,+0.09]$ & yes \\
\midrule
release & PR-AUC(24f)$-$(4s)    & PR-AUC & $+0.549$ & $[+0.46,\,+0.61]$ & yes \\
$1.00\%$ & PR-AUC(24f)$-$(4s) & PR-AUC & $+0.591$ & $[+0.49,\,+0.67]$ & yes \\
$0.50\%$ & PR-AUC(24f)$-$(4s) & PR-AUC & $+0.607$ & $[+0.46,\,+0.74]$ & yes \\
$0.10\%$ & PR-AUC(24f)$-$(4s) & PR-AUC & $+0.216$ & $[-0.05,\,+0.51]$ & no \\
$0.05\%$ & PR-AUC(24f)$-$(4s) & PR-AUC & $+0.334$ & $[+0.10,\,+0.78]$ & yes \\
\bottomrule
\end{tabular}
\end{table}

The third observation is the \textbf{recall@top-1\% monotonicity}: the
24-feature tree recovers $100\%$ of positives in the top-1\% score
slice at every rate down to $0.05\%$ (5 positives). This is the
operating-point evidence that the iter-4 ROC-AUC and iter-8 noise
budget already implied but did not measure: at the deployed operating
point the full tree keeps the alert queue clean enough to recover every
positive at the top-1\% review budget, while the 4-aggregate surrogate
falls to $40\%$ recall at $0.05\%$.

\paragraph{Reading the 4-aggregate surrogate as an oracle LLM sensor.}
A real LLM-as-sensor would face the same positive-rate stress but with
the additional burden of \emph{noisy} numeric outputs. The iter-8
sensor-noise sweep showed that $\sigma \ge 0.05$ on the four
aggregates measurably degrades the 24-feature tree; an LLM sensor that
produces a single deterministic 4-vector per transaction has already
lost the PR-AUC contest by $\approx 0.55$ on the released test split at
the release rate, and the noise budget further restricts it to
$\sigma \le 0.02$ for the tree to register no measurable loss. The
combined picture is a sharp negative-evidence picture for the
LLM-as-sensor pattern on this dataset at any realistic fraud base
rate.

\paragraph{Cross-paper coupling (Pillar 3 / Pillar 4).} The
\secref{sec:p8-scorer} operational stack (latency, cost, calibration,
injection surface) and the \secref{sec:p8-evidence-realistic-ratios}
operating-point gap both reduce to the same diagnostic: an auxiliary
model channel must deliver information strictly orthogonal to the
dominant predictor to register above the noise floor of a paired
bootstrap CI. The Pillar-3 ZVF controller (items 08, 13, 16 of the
P5P8 improvements ledger) reaches the same conclusion in the
RL policy-improvement setting.
